% ch04_functions.tex -- Functions and Scope
\chapter{Functions and Scope}
\label{ch:fonctions}

\begin{intuition}
A function is like a \emph{machine} in a laboratory: you feed it
\emph{inputs} (reagents), it performs a \emph{process} (chemical reaction) and
produces an \emph{output} (product). Rather than rebuilding the machine for each
experiment, you reuse it. This is the \textbf{DRY}\index{DRY} principle:
``\emph{Don't Repeat Yourself}''.
\end{intuition}

% ============================================================
\section{Why use functions?}
\index{function}

\begin{itemize}
    \item \textbf{Reusability}: write a formula once and call it everywhere.
    \item \textbf{Readability}: a program broken into functions is easier to understand.
    \item \textbf{Debugging}: each function can be tested independently.
    \item \textbf{Modularity}\index{modularity}: each function has a single responsibility.
\end{itemize}

% --- TikZ machine ---
\begin{figure}[ht]
\centering
\begin{tikzpicture}[
    box/.style={rectangle, draw, fill=blue!10, rounded corners,
        minimum width=3cm, minimum height=1.5cm, align=center, font=\bfseries},
    io/.style={rectangle, draw, fill=orange!15, rounded corners,
        minimum width=2cm, minimum height=0.8cm, align=center},
    arr/.style={->, thick, >=stealth}
]
\node[io] (in) {Inputs\\(parameters)};
\node[box, right=2cm of in] (fn) {Function\\(processing)};
\node[io, right=2cm of fn] (out) {Output\\(\py{return})};
\draw[arr] (in) -- (fn);
\draw[arr] (fn) -- (out);
\end{tikzpicture}
\caption{A function viewed as a machine: inputs $\to$ processing $\to$ output.}
\label{fig:fonction-machine}
\end{figure}

% ============================================================
\section{Defining and calling a function}
\index{def}\index{return}

\begin{definition}[Function syntax]
\begin{minted}{python}
def function_name(parameter1, parameter2):
    """Docstring: description of the function."""
    # function body
    return result
\end{minted}
\end{definition}

\begin{codeblock}[title=\textbf{Python}]
\begin{minted}{python}
def energie_cinetique(masse, vitesse):
    """Compute the kinetic energy E = 0.5 * m * v^2."""
    return 0.5 * masse * vitesse**2

# Call the function
E = energie_cinetique(2.0, 3.0)
print(f"E = {E} J")
\end{minted}
\end{codeblock}

\begin{codeoutput}
\begin{verbatim}
E = 9.0 J
\end{verbatim}
\end{codeoutput}

% ============================================================
\section{Docstrings}
\index{docstring}

A \emph{docstring} is a documentation string placed right after \py{def}.
It is accessible via \py{help(function_name)}.

\begin{codeblock}[title=\textbf{Python}]
\begin{minted}{python}
def force_gravitationnelle(m1, m2, r):
    """
    Compute the gravitational force between two masses.

    Parameters
    ----------
    m1 : float -- mass 1 (kg)
    m2 : float -- mass 2 (kg)
    r  : float -- distance between centers (m)

    Returns
    -------
    float -- force in Newtons
    """
    G = 6.674e-11  # gravitational constant
    return G * m1 * m2 / r**2

# Earth-Moon force
F = force_gravitationnelle(5.972e24, 7.342e22, 3.844e8)
print(f"Earth-Moon force: {F:.2e} N")
\end{minted}
\end{codeblock}

\begin{codeoutput}
\begin{verbatim}
Earth-Moon force: 1.98e+20 N
\end{verbatim}
\end{codeoutput}

% ============================================================
\section{Default arguments and keyword arguments}
\index{default argument}\index{keyword argument}

\begin{codeblock}[title=\textbf{Python}]
\begin{minted}{python}
def chute_libre(t, g=9.81, v0=0):
    """Position of an object in free fall: y = v0*t + 0.5*g*t^2."""
    return v0 * t + 0.5 * g * t**2

# Using default values
print(f"Earth: y = {chute_libre(2.0):.2f} m")

# On the Moon (g = 1.62 m/s^2)
print(f"Moon : y = {chute_libre(2.0, g=1.62):.2f} m")

# With initial velocity
print(f"With v0=5: y = {chute_libre(2.0, v0=5):.2f} m")
\end{minted}
\end{codeblock}

\begin{codeoutput}
\begin{verbatim}
Earth: y = 19.62 m
Moon : y = 3.24 m
With v0=5: y = 29.62 m
\end{verbatim}
\end{codeoutput}

\begin{bestpractice}
Use keyword arguments to improve readability:
\py{chute\_libre(t=2.0, g=1.62)} is clearer than \py{chute\_libre(2.0, 1.62)}.
\end{bestpractice}

% ============================================================
\section{Functions with multiple return values}

\begin{codeblock}[title=\textbf{Python}]
\begin{minted}{python}
import math

def coordonnees_polaires(x, y):
    """Convert Cartesian coordinates to polar coordinates."""
    r = math.sqrt(x**2 + y**2)
    theta = math.atan2(y, x)
    return r, theta    # returns a tuple

r, angle = coordonnees_polaires(3, 4)
print(f"r = {r:.2f}, theta = {math.degrees(angle):.1f} deg")
\end{minted}
\end{codeblock}

\begin{codeoutput}
\begin{verbatim}
r = 5.00, theta = 53.1 deg
\end{verbatim}
\end{codeoutput}

% ============================================================
\section{Variable scope: local vs global}
\index{scope}\index{local variable}\index{global variable}

\begin{definition}[Scope]
A variable defined \emph{inside} a function is \textbf{local}: it only exists
during the execution of the function. A variable defined outside is
\textbf{global}.
\end{definition}

\begin{codeblock}[title=\textbf{Python}]
\begin{minted}{python}
x = 10          # global variable

def ma_fonction():
    x = 5       # local variable (different!)
    print(f"Inside the function: x = {x}")

ma_fonction()
print(f"Outside: x = {x}")
\end{minted}
\end{codeblock}

\begin{codeoutput}
\begin{verbatim}
Inside the function: x = 5
Outside: x = 10
\end{verbatim}
\end{codeoutput}

\begin{warning}
Avoid using \py{global} to modify global variables from within a
function. \emph{Always} prefer passing values as parameters and using
\py{return}.
\end{warning}

% ============================================================
\section{Lambda functions}
\index{lambda}

\py{lambda} functions are short anonymous functions, useful for simple
operations.

\begin{codeblock}[title=\textbf{Python}]
\begin{minted}{python}
# Lambda function to convert Celsius to Kelvin
celsius_to_kelvin = lambda T: T + 273.15

temperatures_C = [-40, 0, 20, 37, 100]
temperatures_K = [celsius_to_kelvin(T) for T in temperatures_C]

for c, k in zip(temperatures_C, temperatures_K):
    print(f"{c:>4} C = {k:>6.2f} K")
\end{minted}
\end{codeblock}

\begin{codeoutput}
\begin{verbatim}
 -40 C = 233.15 K
   0 C = 273.15 K
  20 C = 293.15 K
  37 C = 310.15 K
 100 C = 373.15 K
\end{verbatim}
\end{codeoutput}

% ============================================================
\section{Common errors}

\begin{errmark}
\begin{enumerate}
    \item \textbf{Forgetting \py{return}}: without \py{return}, a function returns
    \py{None} by default. The computation is performed but the result is lost!
    \begin{minted}{python}
def aire_cercle(r):
    a = 3.14159 * r**2
    # Forgot return a!

resultat = aire_cercle(5)
print(resultat)  # Prints None
    \end{minted}

    \item \textbf{Mutable default argument}: using a list as a default
    value is dangerous, because it is shared across all calls!
    \begin{minted}{python}
# ERROR: the list is shared
def ajouter(element, liste=[]):
    liste.append(element)
    return liste

print(ajouter(1))  # [1]
print(ajouter(2))  # [1, 2] instead of [2]!

# CORRECT: use None
def ajouter(element, liste=None):
    if liste is None:
        liste = []
    liste.append(element)
    return liste
    \end{minted}

    \item \textbf{Confusing \py{print} and \py{return}}: \py{print} displays on screen
    but does not return a value usable in an expression.
\end{enumerate}
\end{errmark}

% ============================================================
\section{Mini-project: Physics Formula Calculator}

\begin{projet}
Create a module containing the following functions:
\begin{itemize}
    \item \py{energie\_cinetique(m, v)}: $E_c = \frac{1}{2}mv^2$
    \item \py{force\_gravitationnelle(m1, m2, r)}: $F = G\frac{m_1 m_2}{r^2}$
    \item \py{periode\_pendule(L, g=9.81)}: $T = 2\pi\sqrt{L/g}$
    \item \py{loi\_ohm(V=None, R=None, I=None)}: $V = RI$ (provide two quantities, compute the third)
\end{itemize}
Each function must have a complete docstring. Test them with known values.
\end{projet}

\begin{codeblock}[title=\textbf{Python}]
\begin{minted}{python}
import math

def periode_pendule(L, g=9.81):
    """Period of a simple pendulum: T = 2*pi*sqrt(L/g)."""
    return 2 * math.pi * math.sqrt(L / g)

def loi_ohm(V=None, R=None, I=None):
    """Ohm's law: V = R * I. Provide two quantities."""
    if V is None:
        return R * I
    elif R is None:
        return V / I
    elif I is None:
        return V / R

# Tests
print(f"Pendulum L=1m: T = {periode_pendule(1.0):.3f} s")
print(f"Ohm's law: V = {loi_ohm(R=100, I=0.5):.1f} V")
print(f"Ohm's law: I = {loi_ohm(V=12, R=100):.3f} A")
\end{minted}
\end{codeblock}

\begin{codeoutput}
\begin{verbatim}
Pendulum L=1m: T = 2.006 s
Ohm's law: V = 50.0 V
Ohm's law: I = 0.120 A
\end{verbatim}
\end{codeoutput}

% ============================================================
\section{Exercises}

\begin{exercise}[$\star$ --- Temperature Conversion]
Write two functions \py{celsius\_to\_fahrenheit(T)} and \py{fahrenheit\_to\_celsius(T)}.
Test with boiling water (100\,\textdegree C = 212\,\textdegree F).
\end{exercise}

\begin{exercise}[$\star$ --- Factorial]
Write a function \py{factorielle(n)} that computes $n!$ using a loop.
Verify that \py{factorielle(10) == 3628800}.
\end{exercise}

\begin{exercise}[$\star\star$ --- Recursive Function]
Rewrite the factorial in a \emph{recursive}\index{recursion} manner: a function that
calls itself. What is the base case?
\end{exercise}

\begin{exercise}[$\star\star$ --- Descriptive Statistics]
Write a function \py{statistiques(donnees)} that receives a list of numbers
and returns the mean, the median, and the standard deviation (without using any library).
\end{exercise}

\begin{exercise}[$\star\star\star$ --- Newton's Method]
Implement Newton's method\index{Newton's method} for finding the zeros
of a function: $x_{n+1} = x_n - \frac{f(x_n)}{f'(x_n)}$.
Your function should take as parameters \py{f}, \py{df} (derivative), \py{x0}
(initial estimate), and \py{tol} (tolerance). Test with
$f(x) = x^2 - 2$ to recover $\sqrt{2}$.
\end{exercise}

% ============================================================
\begin{keyformulas}{Chapter Summary: Functions and Scope}
\begin{itemize}
    \item \textbf{Definition}: \py{def name(params):} ... \py{return result}
    \item \textbf{Docstring}: documentation string in triple quotes
    \item \textbf{Default arguments}: \py{def f(x, n=10):} --- n defaults to 10 if not specified
    \item \textbf{Keyword arguments}: \py{f(x=3, n=20)} for clarity
    \item \textbf{Multiple values}: \py{return a, b} returns a tuple
    \item \textbf{Local scope}: variables inside a function are isolated
    \item \textbf{Lambda}: \py{lambda x: x**2} for short functions
    \item \textbf{DRY principle}: never copy-paste code, write a function
\end{itemize}
\end{keyformulas}
